\documentclass[conference]{IEEEtran}
\IEEEoverridecommandlockouts

\usepackage{cite}
\usepackage{amsmath,amssymb,amsfonts}
\usepackage{graphicx}
\usepackage{textcomp}
\usepackage{xcolor}
\usepackage{booktabs}
\usepackage{array}
\usepackage{url}

\graphicspath{{figures/}{../docs/figures/}}

\begin{document}

\title{Sequential ATPG with Progressive Residual Targeting on\\
Timing-Constrained Partial-Scan Circuits}

\author{
\IEEEauthorblockN{Yen-Chu Lo\textsuperscript{*}, Rueilian Ting\textsuperscript{*}, and Ssu-Wei Huang\textsuperscript{*}}
\IEEEauthorblockA{
Graduate Institute of Electronics Engineering, National Taiwan University\\
EEE5001 VLSI Testing --- Group 5\\
\textsuperscript{*}These authors contributed equally (co-first authors).}
}

\maketitle

\begin{abstract}
In scan-based testing, some flip-flops remain non-scan due to timing
constraints, creating a partial-scan circuit. At $T{=}1$, non-scan flip-flop
(FF) outputs are unknown (X), causing a \textbf{31.26\,pp average gap to
full-scan coverage} across 10 benchmark circuits. We investigate how
multi-frame sequential ATPG can recover this loss. We first establish that
$T{=}2$ sequential ATPG recovers significant coverage, while $T{=}4$ adds
little further benefit, raising the central question of \emph{what mechanism
drives $T{=}2$ recovery}: the frame-based backtrack-limit differential
($T1{=}800$ at $T{=}1$ vs.\ 5000 at $T{>}1$) or Two-Phase State Justification
decoupling. Through five controlled ablation experiments on the same 10
circuits, each varying exactly one parameter, we find: (1) Two-Phase ON raises
average fault coverage (FC) from 73.48\% to \textbf{90.42\%} (+16.94\,pp),
up to +58.49\,pp on individual circuits; (2) uniform $T1{=}5000$ changes FC by
only +0.02\,pp---the backtrack-limit differential is \emph{not} the mechanism;
(3) enhanced backtrace (+0.47\,pp) and static learning (+0.00\,pp) are
negligible; (4) with Two-Phase ON, $T{=}4$ adds 0\,pp universally; (5) the
pipeline narrows the gap to full-scan from 31.26\,pp to \textbf{4.24\,pp}.
\end{abstract}

\begin{IEEEkeywords}
Automatic test pattern generation, partial scan, sequential ATPG, time-frame
expansion, stuck-at faults, state justification.
\end{IEEEkeywords}

\section{Introduction}
\label{sec:intro}
Placeholder --- full conversion of \texttt{docs/final\_report.md} in progress.

\section{Conclusion}
\label{sec:conclusion}
Placeholder.

%% Bibliography enabled after full conversion (refs.bib populated):
%\bibliographystyle{IEEEtran}
%\bibliography{refs}

\end{document}
